\begin{definition}[\textbf{Arborescence}]
	an arborescence is a directed graph where there exists a vertex $r$ (called the root) such that, for any other vertex $v$, there is exactly one directed walk from $r$ to $v$ (noting that the root $r$ is unique). If the underly tree is spanning in \( G_D \), then we say the arborescence is spanning.
\end{definition}

\begin{definition}
	\leavevmode
	\begin{itemize}
		\item \( \cE (D,e) = \) \# of Eulerian tours that has its first edge to be \( e \).
		\item \( \tau (D,v) =  \) \# of (spanning) arborescences rooted at \( v \).
	\end{itemize}
\end{definition}

\begin{definition}[\textbf{Weakly Connected}]
	A digraph \( D \) is \underline{weakly connected} if \( D \) is connected as an undirected graph.
\end{definition}

\begin{definition}[\textbf{Strongly Connected}]
	A digraph \( D \) is \underline{strongly connected} if there is a directed walk between any two vertices. Note that a walk goes from one of them to the other is enough.
\end{definition}

\begin{note}
	Clearly if a graph \( D \) is strongly connected, then it is weakly connected.
\end{note}

\begin{lemma}
	Let \( D \) be a finite directed graph without isolated vertices, \textbf{TFAE}:
	\begin{enumerate}
		\item \( D \) is Eulerian.
		\item \( D \) is connected as an undirected graph, and is balanced. (weakly connected + balanced)
		\item \( D \) is connected directedly and balanced. (strongly connected + balanced)
	\end{enumerate}
	A digraph is balanced if for any vertices \( v \) in \( D \), the indegree of \( v \) is equal to the outdegree of \( v \) is equal to the outdegree of \( v \).
\end{lemma}

\begin{proof}
	Suppose that \( D \) is eulerian and there exists an eulerian tour \( e_1, \ldots, e_q \), with
	the source being \( v \). As we move along the tour, whenever we enter a vertex we must then exit
	it, except the final vertex of \( e_f \) without exiting it. But this is the source vertex so we
	exit it at the beginning, thus every vertex entered as often as exit, and they must have the same
	indegree and outdegree. Therefore \( D \) is balanced. Note that \( D \) is clearly strongly
	connected and thus weakly connected.

	Conversely given a balanced, weakly connected digraph, we want to see that it is eulerian.
	The generalhidea is to first prove that for any vertex \( v \) in the digraph, there exists a tour
	starts with \( v \) and then ends with \( v \). And then since \( D \) is a finite digraph, one
	can then choose a tour with maximum length \( r \), and try to prove that \( r=q \) be proceeding
	contradiction that if \( r < q \), such will contradict to the maximality of \( r \).

	Since \( D \) is connected one may assume that \( D \) contains at least one edge \( e \), take
	the source of that edge and denote it as \( v \). Consider \( e = (v,u) \), if \( u = v\) and we
	are done. If \( u\ne v \), since the indegree of \( u \) equals to the outdegree, then there is
	some exit edge \( e_1 \), either the target of \( e_1 \) being \( v \) and we are done, or being
	some vertex \( w \), and continuing this way, either we complete a tour or else we have entered
	the current vertex once more than we have exited, which guaranteed by balanced that there exists
	an exit along a new edge other than previously in the tour. Since \( D \) has finitely many edges,
	evetually we must complete a tour. Thus \( D \) does have a tour that starts at any
	vertex \( v \), and thus for any edge \( e \), there exists a tour that use \( e \).

	One can then take the maximum size of tour, denoted as \( e_1, \ldots, e_r \), one shall then want
	to prove that \( r = q \). Suppose that \( r < q \), then the edges form a subgrpah \( C \) of \( D \).
	By definition, \( C \) should be also a balanced graph. Thus consider the graph
	\( H:= D \backslash C \) by removing all edges of \( C \) out of \( D \), since \( D,C \) are both
	balanced graph, \( H \) shall also be a balanced graph, although may not connected. One shall then
	observe that there must exists a connected component \( H_1 \) in \( H \) that contains at least
	one edge, and a vertex \( v \) that is also contained in \( C \), otherwise see that the tour is
	then directly being the eulerian tour and we have \( r = q \), a contradiction. Thus we have such
	\( H_1 \) exists. But then see that in \( H_1 \) there exists a tour \( C' \) starts from \( v \)
	and ends with \( v \) by our previous reasoning, this gives us a ``detour'' as one can enlarge the
	tour by when moving along \( C \), when arrive \( v \), one can make detour on \( C' \) and back
	on \( v \) then proceed. This gives us a longer tour, contradicting to the maximality of \( r \)
	$\lightning$. Thus \( r = q \) and yields the existence of Eulerian tour, and thus \( D \) is
	Eulerian.
\end{proof}

\begin{theorem}[\textbf{BEST Theorem | de Bruijn, van Aardenne-Ehrenfest, Smith, Tutte Theorem}]
	Let \( D \) be a connected balanced digraph with vertex set being \( V \). Fix an edge \( e \) of \( D \) and let \( v = init(e) \). Then
	\[
		\cE(D,e) = \tau(D,v) \prod_{u\in V} \bace{\outdeg_G(u)-1}!
	\]
\end{theorem}

\begin{proof}

\end{proof}
